\section{Calculabilité}\label{section-computability}

\begin{center}
\textcolor{red}{\textbf{Premier brouillon produit avec l'aide d'un
modèle d'IA.  Ce texte doit impérativement être relu, corrigé et
vérifié par l'enseignant avant toute diffusion.}}
\end{center}

\subsection{Introduction}\label{subsection-computability-introduction}

\thingy La \defin{calculabilité} se situe à l'interface entre la logique
mathématique et l'informatique théorique.  Elle est née de
préoccupations venues de la logique, notamment des travaux de Hilbert
et de Gödel, et elle a conduit à certains des premiers concepts
fondamentaux de l'informatique théorique, comme le $\lambda$-calcul et
la machine de Turing.

Son but est d'étudier les limites de ce que peut ou ne peut pas faire
un algorithme.  Contrairement à la théorie de la complexité, elle ne
cherche pas à mesurer précisément le temps ou la mémoire nécessaires à
un calcul.  Elle autorise des ressources arbitrairement grandes, à
condition qu'elles restent finies pour chaque calcul qui termine.  Elle
permet ainsi de démontrer des résultats négatifs de la forme
«~\emph{aucun} algorithme ne peut résoudre ce problème~».

On peut également considérer des notions relatives de calculabilité,
dans lesquelles un algorithme dispose d'informations supplémentaires
sous la forme d'un oracle, ainsi que diverses généralisations dites de
calculabilité supérieure.  Ces développements ne seront toutefois pas
au centre de cette partie.

\thingy Quelques noms permettent de situer historiquement le sujet.
Le mot «~algorithme~» dérive du nom de Muḥammad ibn Mūsá
al-\b{H}wārizmī (vers 780--vers 850).  Blaise Pascal
(1623--1662) construisit une machine à calculer, tandis que Charles
Babbage (1791--1871) conçut son \textit{Analytical Engine}, machine
programmable qui, dans une idéalisation convenable, aurait eu la
puissance de calcul d'une machine de Turing.  Ada, née Byron, comtesse
de Lovelace (1815--1852), est notamment associée aux premiers travaux
de programmation pour cette machine.

Du côté mathématique, Richard Dedekind (1831--1916) introduisit des
définitions par récursion primitive.  David Hilbert (1862--1943)
formula l'\textit{Entscheidungsproblem}, qui demandait essentiellement
une procédure générale permettant de décider la vérité des énoncés
mathématiques d'une certaine forme.  Jacques Herbrand (1908--1931),
Kurt Gödel (1906--1978) et Stephen C. Kleene (1909--1994)
contribuèrent à la théorie des fonctions récursives.  Alonzo Church
(1903--1995) introduisit le $\lambda$-calcul, tandis qu'Alan M. Turing
(1912--1954) définit les machines qui portent son nom et démontra
l'indécidabilité du problème de l'arrêt.  Emil Post (1897--1954)
étudia notamment les ensembles calculablement énumérables.  Enfin,
Haskell Curry (1900--1982) développa la logique combinatoire et
contribua à mettre en évidence les liens entre preuves et typage.

\thingy Une première tentative, encore informelle, consiste à dire
qu'une fonction $f$ est \defin{calculable} lorsqu'il existe un
algorithme qui, prenant en entrée un élément $x$ du domaine de
définition de $f$, termine en temps fini et renvoie la valeur $f(x)$.

Cette formulation fait immédiatement apparaître plusieurs difficultés.
Que faut-il entendre exactement par «~algorithme~»~?  Quelles sortes
de valeurs un algorithme peut-il accepter ou renvoyer~?  Comment
décrire les situations dans lesquelles il ne termine pas~?  Enfin, il
faut distinguer le programme lui-même, ou son \emph{intention}, de la
fonction qu'il calcule, ou son \emph{extension}~: plusieurs programmes
différents peuvent calculer exactement la même fonction.

\thingy La calculabilité ne s'intéresse pas à l'efficacité des
algorithmes, mais seulement à leur terminaison en temps fini.  Par
exemple, pour déterminer si un entier $n$ est premier, on peut tester
tous les produits $i\times j$ avec $2\leq i,j\leq n-1$.  Cet
algorithme est extrêmement inefficace, mais il termine pour toute
entrée $n$ et suffit donc du point de vue de la calculabilité.

De même, la calculabilité n'a pas peur des grands entiers.  Considérons
la \defin{fonction d'Ackermann} définie par
\[
\begin{aligned}
A(m,n,0) &= m+n,\\
A(m,1,k+1) &= m,\\
A(m,n+1,k+1) &= A(m,\,A(m,n,k+1),\,k).
\end{aligned}
\]
Cette définition décrit un algorithme par appels récursifs, et ces
appels terminent.  La fonction est donc calculable.  Ses valeurs
deviennent néanmoins gigantesques très rapidement~:
\[
A(2,6,3)
 = 2^{2^{2^{2^{2^2}}}}
 = 2^{2^{65\,536}},
\]
et $A(2,4,4)=A(2,65\,536,3)$ est déjà inimaginablement grand.
Une valeur telle que $A(100,100,100)$ est tout à fait hors de portée
d'un ordinateur réel, sans que cela empêche la fonction $A$ d'être
calculable.

\thingy On peut aborder la notion de calculabilité de plusieurs
manières.

L'approche informelle considère un algorithme comme un calcul
\emph{finitiste}, mené par un humain ou une machine selon des
instructions précises, en temps fini et sur des données finies.

L'approche pragmatique identifie les fonctions calculables à celles
que l'on peut programmer dans un langage général dit
«~Turing-complet~», par exemple Python, Java, C ou OCaml, à condition
d'idéaliser ce langage.  On suppose notamment que les entiers n'ont
aucune limite de taille et que la mémoire disponible n'est pas bornée
à l'avance.  On écarte également les véritables sources de hasard et
de non-déterminisme.

Enfin, plusieurs approches formelles ont été proposées~:
\begin{itemize}
\item les fonctions générales récursives de Herbrand, Gödel et Kleene~;
\item le $\lambda$-calcul de Church, proche des langages fonctionnels~;
\item les machines de Turing~;
\item les machines à registres, notamment étudiées par Post.
\end{itemize}

\thingy Le résultat mathématique fondamental est que les principales
formalisations précédentes conduisent à la même classe de fonctions.

\begin{thm}[Post, Turing]\label{church-turing-equivalence}
Les classes de fonctions partielles
$\mathbb{N}\dasharrow\mathbb{N}$ suivantes coïncident~:
\begin{enumerate}
\item les fonctions générales récursives~;
\item les fonctions représentables dans le $\lambda$-calcul~;
\item les fonctions calculables par une machine de Turing.
\end{enumerate}
\end{thm}

On parle donc de \defin{calculabilité au sens de Church-Turing}.
Tous les langages de programmation généraux usuels, convenablement
idéalisés, calculent eux aussi exactement ces fonctions~: c'est ce que
signifie, du point de vue de la calculabilité, l'expression
«~Turing-complet~».

Il faut distinguer ce théorème de deux affirmations plus générales.  La
\defin[Church-Turing (thèse de)]{thèse de Church-Turing}, de nature
philosophique, affirme que cette définition formelle saisit exactement
la notion intuitive d'algorithme finitiste.  Une conjecture physique
plus forte affirme que les calculs réalisables mécaniquement dans
l'Univers, avec un temps et une énergie finis mais sans borne fixée à
l'avance, sont précisément ceux de Church-Turing.  Cette conjecture
reste la même si l'on autorise un ordinateur quantique~: un ordinateur
quantique pourrait apporter des gains de complexité, mais ne devrait
pas calculer davantage de fonctions au sens considéré ici.

\thingy Nous allons étudier trois descriptions de ces mêmes fonctions
calculables.

Les \defin{fonctions générales récursives} sont particulièrement
commodes du point de vue mathématique.  Toutes les données y sont
représentées par des entiers, les fonctions considérées étant de la
forme
\[
\mathbb{N}^k\dasharrow\mathbb{N}.
\]
Leur définition est inductive et permet d'associer des numéros aux
programmes.

Les \defin{machines de Turing} constituent une modélisation très simple
d'un ordinateur.  Elles travaillent sur une bande illimitée \textit{a
priori}, qui joue le rôle de mémoire.  Leur caractère algorithmique est
particulièrement manifeste.  Elles sont aussi commodes pour étudier la
complexité, même si cette question ne sera pas notre préoccupation
principale ici.

Enfin, le \defin{$\lambda$-calcul} pur non typé est un système
symbolique proche des langages de programmation fonctionnels comme
Lisp, Haskell ou OCaml.  Il est relativement facile d'y écrire de
véritables programmes, mais son étude fait intervenir plusieurs
subtilités liées notamment aux variables et à la stratégie
d'évaluation.

\subsection{Données finies, codages et fonctions partielles}
\label{subsection-finite-data-and-partial-functions}

\thingy Un algorithme travaille sur des \defin{données finies}.  Cette
expression recouvre les booléens, les entiers, les chaînes de
caractères, les listes, les tableaux et les structures finies plus
complexes, par exemple les graphes finis.

Deux attitudes sont possibles.  Le \defin{typage} distingue les
différentes sortes de données et précise les opérations permises sur
chacune d'elles.  Cette approche est élégante et proche de
l'informatique réelle.  À l'opposé, le \defin[Gödel
(codage de)]{codage de Gödel} représente toutes les données finies par
des entiers naturels.  Cette seconde approche est moins élégante mais
simplifie considérablement la définition de la calculabilité, puisqu'il
suffit alors d'étudier des fonctions
$\mathbb{N}^k\dasharrow\mathbb{N}$, voire seulement
$\mathbb{N}\dasharrow\mathbb{N}$.

\thingy\label{godel-coding} Pour coder un couple d'entiers, on peut par
exemple poser
\[
\langle m,n\rangle
 := m+\frac{1}{2}(m+n)(m+n+1).
\]
Cette formule définit une bijection calculable
$\mathbb{N}^2\to\mathbb{N}$, dont la réciproque est elle aussi
calculable.  On peut donc aussi bien manipuler un couple $(m,n)$ que
l'entier $\langle m,n\rangle$ qui le représente.

On peut ensuite coder une liste finie par
\[
\dbllangle a_0,\ldots,a_{k-1}\dblrangle
 := \langle a_0,
       \langle a_1,
       \langle\cdots,
       \langle a_{k-1},0\rangle+1
       \cdots\rangle+1\rangle+1,
\]
avec la convention $\dbllangle\dblrangle:=0$.  On obtient ainsi une
bijection calculable entre l'ensemble des suites finies d'entiers
naturels et $\mathbb{N}$.

Des codages analogues permettront de représenter des expressions, des
arbres, des configurations de machines et même des programmes par des
entiers.  Le choix précis du codage est sans importance pour la théorie,
pourvu que les opérations d'encodage et de décodage nécessaires soient
calculables.

{\footnotesize
\thingy \emph{Attention~:} ces codages sont des outils théoriques.  Ils
ne constituent généralement pas de bonnes représentations pour un
programme réel, en raison notamment de la taille considérable des
entiers obtenus.
\par}

\thingy Même si l'on cherche principalement à décrire les algorithmes
qui terminent, la définition générale de la calculabilité doit prendre
en compte ceux qui ne terminent pas.  On est ainsi conduit à travailler
avec des \defin{fonctions partielles}.

Une fonction partielle
\[
f\colon\mathbb{N}\dasharrow\mathbb{N}
\]
est une fonction $D\to\mathbb{N}$ définie sur une partie
$D\subseteq\mathbb{N}$.  On écrit
\[
f(n)\downarrow
\]
lorsque $n\in D$, c'est-à-dire lorsque $f(n)$ est définie, et
\[
f(n)\uparrow
\]
lorsque $n\notin D$.

L'expression
\[
f(n)\downarrow=g(m)
\]
signifie que $f(n)$ et $g(m)$ sont toutes deux définies et ont la même
valeur.  Par convention, l'égalité
\[
f(n)=g(m)
\]
entre deux expressions éventuellement partielles signifie qu'elles
sont définies simultanément et qu'elles ont alors la même valeur.
Certains auteurs écrivent plutôt $f(n)\simeq g(m)$ pour cette relation.

Enfin, si l'un des arguments $g_i(\underline{x})$ n'est pas défini, on
convient que
\[
h(g_1(\underline{x}),\ldots,g_k(\underline{x}))\uparrow.
\]
Une fonction partielle définie sur tout $\mathbb{N}^k$ est dite
\defin[totale (fonction)]{totale}.  Une fonction totale est donc un cas
particulier de fonction partielle.

\begin{defn}\label{definition-computable-partial-function}
Une fonction partielle
$f\colon\mathbb{N}^k\dasharrow\mathbb{N}$ est
\defin[calculable (fonction)]{calculable} lorsqu'il existe un algorithme
qui, pour toute entrée $\underline{x}\in\mathbb{N}^k$,
\begin{itemize}
\item termine en temps fini et renvoie $f(\underline{x})$ lorsque
  $f(\underline{x})\downarrow$~;
\item ne termine pas lorsque $f(\underline{x})\uparrow$.
\end{itemize}
\end{defn}

\begin{defn}\label{definition-decidable-semidecidable}
Une partie $A\subseteq\mathbb{N}^k$ est dite
\defin[décidable (ensemble)]{décidable} lorsque sa fonction indicatrice
\[
\mathbf{1}_A\colon\underline{x}\mapsto
\begin{cases}
1&\text{si $\underline{x}\in A$,}\\
0&\text{si $\underline{x}\notin A$}
\end{cases}
\]
est calculable.

Elle est dite \defin[semi-décidable
(ensemble)]{semi-décidable} lorsque la fonction partielle
\[
\underline{x}\mapsto
\begin{cases}
1&\text{si $\underline{x}\in A$,}\\
\uparrow&\text{si $\underline{x}\notin A$}
\end{cases}
\]
est calculable.
\end{defn}

Ainsi, un algorithme de décision répond toujours «~oui~» ou «~non~».
Un algorithme de semi-décision répond «~oui~» lorsque l'entrée
appartient à $A$, mais peut calculer indéfiniment lorsqu'elle n'y
appartient pas.

\thingy Le mot «~récursif~» possède plusieurs sens apparentés mais
distincts.  Il peut signifier «~défini par récurrence~», comme dans les
expressions «~fonction primitive récursive~» et «~fonction générale
récursive~».  Il est aussi employé comme synonyme historique de
«~calculable~».  Enfin, en informatique, un appel est dit récursif
lorsqu'une fonction fait appel à elle-même.  Dans la suite, nous
réserverons autant que possible l'expression «~appel récursif~» à ce
dernier sens.

\subsection{Fonctions primitives récursives}
\label{subsection-primitive-recursive-functions}

\thingy Avant de définir les fonctions générales récursives, qui
coïncideront avec les fonctions calculables, commençons par une classe
plus restreinte~: les \defin[primitive récursive
(fonction)]{fonctions primitives récursives}, abrégées en
\textbf{p.r.}

Cette classe est historiquement antérieure à la définition de la
calculabilité au sens de Church-Turing.  Elle est utile comme première
approximation et correspond informatiquement à des programmes dont
toutes les boucles sont bornées \textit{a priori}.  Une très grande
partie des algorithmes usuels définit des fonctions primitives
récursives, mais ce n'est pas le cas de tous les algorithmes
terminants~: la fonction d'Ackermann fournira un contre-exemple.

\begin{defn}\label{primitive-recursive-definition}
La classe $\mathbf{PR}$ des fonctions primitives récursives est la plus
petite classe de fonctions
$\mathbb{N}^k\dasharrow\mathbb{N}$, pour $k$ variable, qui satisfait
les propriétés suivantes.
\begin{itemize}
\item Elle contient les projections
\[
(x_1,\ldots,x_k)\mapsto x_i.
\]
\item Elle contient toutes les fonctions constantes
\[
\underline{x}\mapsto c.
\]
\item Elle contient la fonction successeur $x\mapsto x+1$.
\item Elle est stable par composition~: si
\[
g_1,\ldots,g_\ell\colon\mathbb{N}^k\dasharrow\mathbb{N}
\quad\text{et}\quad
h\colon\mathbb{N}^\ell\dasharrow\mathbb{N}
\]
sont primitives récursives, alors
\[
\underline{x}\mapsto
h(g_1(\underline{x}),\ldots,g_\ell(\underline{x}))
\]
est primitive récursive.
\item Elle est stable par \defin{récursion primitive}~: si
\[
g\colon\mathbb{N}^k\dasharrow\mathbb{N}
\quad\text{et}\quad
h\colon\mathbb{N}^{k+2}\dasharrow\mathbb{N}
\]
sont primitives récursives, alors la fonction
$f\colon\mathbb{N}^{k+1}\dasharrow\mathbb{N}$ définie par
\[
\begin{aligned}
f(\underline{x},0) &= g(\underline{x}),\\
f(\underline{x},z+1)
 &=h(\underline{x},f(\underline{x},z),z)
\end{aligned}
\]
est primitive récursive.
\end{itemize}
\end{defn}

Toutes les fonctions primitives récursives sont en fait totales.  La
notation $\dasharrow$ reste néanmoins commode, notamment lorsque nous
coderons les programmes~: certains entiers ne coderont aucun programme
bien formé et correspondront alors à une fonction non définie.

\thingy L'addition est primitive récursive, car la fonction
$f(x,z)=x+z$ vérifie
\[
\begin{aligned}
f(x,0)&=x,\\
f(x,z+1)&=f(x,z)+1.
\end{aligned}
\]
La fonction initiale $x\mapsto x$ est une projection, et la fonction
$(x,y,z)\mapsto y+1$ s'obtient par composition de la projection sur
$y$ et du successeur.

La multiplication est ensuite primitive récursive, puisque
$f(x,z)=x\cdot z$ vérifie
\[
\begin{aligned}
f(x,0)&=0,\\
f(x,z+1)&=f(x,z)+x.
\end{aligned}
\]
En utilisant de nouveau une récursion primitive, on obtient que
$(x,z)\mapsto x^z$ est primitive récursive.

La fonction définie par
\[
f(x,y,z)=
\begin{cases}
x&\text{si $z=0$,}\\
y&\text{si $z\geq1$}
\end{cases}
\]
est elle aussi primitive récursive, car
\[
\begin{aligned}
f(x,y,0)&=x,\\
f(x,y,z+1)&=y.
\end{aligned}
\]

Enfin, la fonction prédécesseur tronqué
\[
u\mapsto\max(u-1,0)
\]
est primitive récursive.  On peut en déduire que le sont également la
soustraction tronquée
\[
(u,v)\mapsto\max(u-v,0),
\]
le reste $(u,v)\mapsto u\mathbin{\%}v$ et le quotient entier
$(u,v)\mapsto\lfloor u/v\rfloor$.

\review{Les diapositives présentent la construction explicite du
prédécesseur tronqué et des opérations de division comme exercice.
La démonstration détaillée pourra être ajoutée lors de l'intégration
des feuilles d'exercices.}

\thingy Une autre façon de comprendre la classe $\mathbf{PR}$ consiste
à employer un langage de programmation idéalisé dans lequel~:
\begin{itemize}
\item les variables contiennent des entiers naturels de taille
  arbitraire~;
\item les constantes, affectations, tests d'égalité et conditionnelles
  sont disponibles~;
\item les opérations arithmétiques élémentaires sont disponibles~;
\item les appels de fonctions sont permis, mais pas les appels
  récursifs~;
\item toutes les boucles ont un nombre d'itérations borné
  \textit{a priori}.
\end{itemize}
Les programmes écrits dans un tel langage terminent par construction.
Les fonctions qu'ils calculent sont exactement les fonctions
primitives récursives.

Les fonctions de codage et de décodage des couples introduites
en~\ref{godel-coding}, à savoir
\[
(m,n)\mapsto\langle m,n\rangle,\qquad
\langle m,n\rangle\mapsto m,\qquad
\langle m,n\rangle\mapsto n,
\]
sont elles aussi primitives récursives.

\thingy La classe $\mathbf{PR}$ peut aussi être interprétée comme une
classe de complexité extrêmement large.  En anticipant sur la
définition des machines de Turing, la fonction qui, à une machine $M$
et une configuration $C$, associe la configuration suivante $C'$ est
primitive récursive.  Par récursion primitive, la fonction
\[
(n,M,C)\mapsto C^{(n)}
\]
qui associe à $M$, $C$ et $n$ la configuration atteinte après $n$
étapes est donc primitive récursive.

Il s'ensuit qu'une fonction calculée par une machine de Turing en un
nombre d'étapes borné par une fonction primitive récursive est
elle-même primitive récursive.  Réciproquement, toute fonction
primitive récursive est calculable avec une borne primitive récursive
sur sa complexité.

\review{L'équivalence précise entre «~fonction primitive récursive~»
et «~fonction de complexité primitive récursive~» dépend du modèle de
calcul et des conventions de codage.  La formulation et la preuve
devront être vérifiées lors de la rédaction de la partie consacrée aux
machines de Turing.}

En particulier, on a l'inclusion
\[
\mathbf{EXPTIME}\subseteq\mathbf{PR}.
\]

\thingy La fonction d'Ackermann montre néanmoins que toutes les
fonctions calculables totales ne sont pas primitives récursives.  Pour
$m=2$, elle vérifie
\[
\begin{aligned}
A(2,n,0)&=2+n,\\
A(2,1,k+1)&=2,\\
A(2,n+1,k+1)&=A(2,\,A(2,n,k+1),\,k).
\end{aligned}
\]
Cette définition n'est pas une récursion primitive~: dans la dernière
ligne, la valeur de $k$ diminue lors de l'appel extérieur, mais un
appel imbriqué à $A$ intervient dans le deuxième argument.  Il ne
s'agit donc pas d'une simple récursion primitive sur une variable.

On peut montrer que, pour toute fonction primitive récursive
$f\colon\mathbb{N}^k\to\mathbb{N}$, il existe un entier $r=r(f)$ tel
que
\[
f(x_1,\ldots,x_k)
 \leq A(2,x_1+\cdots+x_k+3,r).
\]
Par conséquent, la fonction diagonale
\[
r\mapsto A(2,r,r)
\]
n'est pas primitive récursive.  Elle est pourtant définie par un
algorithme clair qui termine pour toute entrée.

\review{La démonstration de la majoration des fonctions primitives
récursives par les niveaux de la fonction d'Ackermann n'est pas donnée
dans les diapositives.  Elle est seulement annoncée et n'est donc pas
développée ici.}

\subsection{Numérotation des fonctions primitives récursives}
\label{subsection-numbering-primitive-recursive-functions}

\thingy Nous voulons maintenant coder les programmes primitifs
récursifs par des entiers.  Pour certains entiers $e\in\mathbb{N}$,
nous définirons une fonction
\[
\psi_e^{(k)}\colon\mathbb{N}^k\dasharrow\mathbb{N},
\]
appelée la fonction d'arité $k$ \defin[codée
(fonction)]{codée} par $e$.  L'entier $e$ doit être compris comme un
code source ou un programme.

Toute fonction primitive récursive
$f\colon\mathbb{N}^k\to\mathbb{N}$ sera égale à
$\psi_e^{(k)}$ pour au moins un entier $e$.  Cet entier décrit la
manière dont $f$ est construite à l'aide des clauses de la
définition~\ref{primitive-recursive-definition}.  Il n'est pas unique~:
une même fonction peut être calculée par de nombreux programmes
distincts,
\[
f=\psi_{e_1}^{(k)}
 =\psi_{e_2}^{(k)}
 =\cdots.
\]
L'indice $e$ décrit l'intention, tandis que la fonction $f$ est
l'extension.

\thingy Avec le codage de listes de~\ref{godel-coding}, on définit
$\psi_e^{(k)}$ par induction sur la construction primitive récursive.
\begin{itemize}
\item Si $e=\dbllangle0,k,i\dblrangle$, alors
\[
\psi_e^{(k)}(x_1,\ldots,x_k)=x_i.
\]
\item Si $e=\dbllangle1,k,c\dblrangle$, alors
\[
\psi_e^{(k)}(x_1,\ldots,x_k)=c.
\]
\item Si $e=\dbllangle2\dblrangle$, alors
\[
\psi_e^{(1)}(x)=x+1.
\]
\item Si
\[
e=\dbllangle3,k,d,c_1,\ldots,c_\ell\dblrangle,
\qquad
g_i:=\psi_{c_i}^{(k)},\qquad h:=\psi_d^{(\ell)},
\]
alors
\[
\psi_e^{(k)}(\underline{x})
 =h(g_1(\underline{x}),\ldots,g_\ell(\underline{x})).
\]
\item Si $e=\dbllangle4,k,d,c\dblrangle$, avec
$g:=\psi_c^{(k)}$ et $h:=\psi_d^{(k+2)}$, alors
\[
\begin{aligned}
\psi_e^{(k+1)}(\underline{x},0)
 &=g(\underline{x}),\\
\psi_e^{(k+1)}(\underline{x},z+1)
 &=h(\underline{x},
     \psi_e^{(k+1)}(\underline{x},z),z).
\end{aligned}
\]
\end{itemize}
Dans tous les autres cas, $\psi_e^{(k)}$ n'est pas définie.

Par définition, une fonction
$f\colon\mathbb{N}^k\dasharrow\mathbb{N}$ est primitive récursive si
et seulement s'il existe $e\in\mathbb{N}$ tel que
\[
f=\psi_e^{(k)}.
\]

\thingy La numérotation précédente permet de manipuler
algorithmiquement les programmes.  L'opération fondamentale suivante
consiste à fixer certains de leurs arguments.

\begin{thm}[Théorème s-m-n, version primitive récursive]
\label{smn-primitive-recursive}
Pour tous $m,n\in\mathbb{N}$, il existe une fonction primitive
récursive
\[
s_{m,n}\colon\mathbb{N}^{m+1}\to\mathbb{N}
\]
telle que
\[
\psi^{(n)}_{s_{m,n}(e,x_1,\ldots,x_m)}
       (y_1,\ldots,y_n)
 =
\psi^{(m+n)}_e
       (x_1,\ldots,x_m,y_1,\ldots,y_n)
\]
pour tous $e,\underline{x},\underline{y}$.
\end{thm}

\begin{proof}
Avec les conventions de codage précédentes, on peut prendre
\[
\begin{aligned}
s_{m,n}(e,\underline{x})
 :=\dbllangle
 &3,n,e,
 \dbllangle1,n,x_1\dblrangle,\ldots,
 \dbllangle1,n,x_m\dblrangle,\\
 &\dbllangle0,n,1\dblrangle,\ldots,
 \dbllangle0,n,n\dblrangle
 \dblrangle.
\end{aligned}
\]
Ce code décrit la composition du programme $e$ avec, pour ses $m$
premiers arguments, les fonctions constantes de valeurs
$x_1,\ldots,x_m$, et, pour ses $n$ derniers arguments, les projections.
Toutes les opérations de codage utilisées sont primitives récursives,
donc $s_{m,n}$ l'est aussi.
\end{proof}

En clair, $s_{m,n}$ transforme un programme prenant $m+n$ arguments en
un programme dans lequel les $m$ premiers arguments ont été fixés aux
valeurs $x_1,\ldots,x_m$, tandis que les $n$ derniers restent variables.

\subsection{Auto-référence et théorème de récursion}
\label{subsection-kleene-recursion-theorem-pr}

\thingy L'\defin[Quine (astuce de)]{astuce de Quine} est une technique
permettant à un programme d'avoir accès à son propre code source.  Son
idée peut être illustrée par la phrase autoréférente suivante~:

\begin{center}
\emph{Les mots suivants suivis des mêmes mots entre guillemets forment
une phrase intéressante~: «~les mots suivants suivis des mêmes mots
entre guillemets forment une phrase intéressante~».}
\end{center}

La première partie de la phrase sert de modèle et la seconde lui
fournit une représentation d'elle-même.  Une construction analogue
permet d'écrire un programme qui transmet son propre code source à une
fonction arbitraire.  Ainsi, même si un langage ne fournit aucune
instruction primitive donnant accès au code du programme courant, il
est possible de simuler cet accès.

\review{Le pseudocode de la diapositive est volontairement informel.
Il pourra être réintroduit sous forme d'exercice lors de l'intégration
des feuilles d'exercices.}

\begin{thm}[Kleene, version primitive récursive]
\label{kleene-recursion-theorem-pr}
Si $h\colon\mathbb{N}^{k+1}\dasharrow\mathbb{N}$ est primitive
récursive, il existe un entier $e$ tel que
\[
(\forall\underline{x})\qquad
\psi_e^{(k)}(\underline{x})=h(e,\underline{x}).
\]
Plus précisément, il existe une fonction primitive récursive
$b\colon\mathbb{N}^2\to\mathbb{N}$ telle que, si $e:=b(k,d)$, alors
\[
(\forall d)(\forall\underline{x})\qquad
\psi_e^{(k)}(\underline{x})
 =\psi_d^{(k+1)}(e,\underline{x}).
\]
\end{thm}

\begin{proof}
Soit $s:=s_{1,k}$ la fonction fournie par le théorème s-m-n.  La
fonction
\[
(t,\underline{x})\mapsto h(s(t,t),\underline{x})
\]
est primitive récursive.  Il existe donc un entier $c$ tel que
\[
\psi_c^{(k+1)}(t,\underline{x})
 =h(s(t,t),\underline{x}).
\]
Posons $e:=s(c,c)$.  Le théorème s-m-n donne alors
\[
\begin{aligned}
\psi_e^{(k)}(\underline{x})
 &=\psi_{s(c,c)}^{(k)}(\underline{x})\\
 &=\psi_c^{(k+1)}(c,\underline{x})\\
 &=h(s(c,c),\underline{x})\\
 &=h(e,\underline{x}).
\end{aligned}
\]
Les constructions qui, à $d$, associent un indice $c$, puis à $c$
l'indice $s(c,c)$, peuvent être effectuées primitivement récursivement,
ce qui donne la fonction uniforme $b$ annoncée.
\end{proof}

La moralité est que l'on peut donner à un programme accès à son propre
numéro, donc à une représentation de son code source, sans augmenter
la classe des fonctions qu'il peut calculer.

\thingy Il n'existe cependant pas d'interpréteur universel
\emph{primitif récursif} pour les programmes primitifs récursifs.

\begin{thm}\label{primitive-recursive-no-universality}
Il n'existe pas de fonction primitive récursive
$u\colon\mathbb{N}^2\to\mathbb{N}$ telle que
\[
u(e,x)=\psi_e^{(1)}(x)
\]
chaque fois que $\psi_e^{(1)}(x)\downarrow$.
\end{thm}

\begin{proof}
Supposons par l'absurde qu'une telle fonction $u$ existe.  La fonction
\[
(e,x)\mapsto u(e,x)+1
\]
est alors primitive récursive.  Par le théorème de récursion de Kleene,
il existe un indice $e$ tel que
\[
\psi_e^{(1)}(x)=u(e,x)+1
\]
pour tout $x$.  Mais la propriété supposée de $u$ donne simultanément
\[
u(e,x)=\psi_e^{(1)}(x),
\]
d'où
\[
u(e,x)=u(e,x)+1,
\]
ce qui est impossible.
\end{proof}

Autrement dit, un interpréteur du langage primitif récursif ne peut pas
être lui-même primitif récursif.  C'est un argument diagonal, analogue
à celui de Cantor.  Il repose à la fois sur le théorème s-m-n et sur le
fait que les fonctions primitives récursives sont totales.

On peut également retrouver cette absence d'universalité à l'aide de
la fonction d'Ackermann.  Pour chaque $k$ fixé, la fonction
\[
(m,n)\mapsto A(m,n,k)
\]
est primitive récursive, et on peut calculer primitivement
récursivement, en fonction de $k$, un indice $a(k)$ tel que
\[
\psi_{a(k)}^{(2)}(m,n)=A(m,n,k).
\]
Si la fonction universelle $(e,n)\mapsto\psi_e^{(1)}(n)$ était
primitive récursive, le théorème s-m-n permettrait d'en déduire que
$(n,k)\mapsto A(2,n,k)$ est primitive récursive, ce qui contredit les
limitations établies précédemment.

\subsection{Fonctions générales récursives}
\label{subsection-general-recursive-functions}

\thingy Les fonctions primitives récursives ne couvrent pas tous les
algorithmes.  Les programmes primitifs récursifs terminent toujours,
n'autorisent que des boucles bornées à l'avance et ne peuvent pas
s'interpréter eux-mêmes.  Pour obtenir une classe correspondant à la
calculabilité générale, il faut autoriser des recherches non bornées,
donc accepter que certains calculs ne terminent pas.

\begin{defn}\label{definition-mu-operator}
Soit
\[
g\colon\mathbb{N}^{k+1}\dasharrow\mathbb{N}.
\]
La fonction $\mu g\colon\mathbb{N}^k\dasharrow\mathbb{N}$ est définie
en posant que $\mu g(\underline{x})$ est le plus petit entier $z$ tel
que
\[
g(z,\underline{x})=0
\]
et tel que $g(i,\underline{x})\downarrow$ pour tout $0\leq i<z$, s'il
existe un tel entier.

Autrement dit, $\mu g(\underline{x})=z$ signifie
\[
g(z,\underline{x})=0
\quad\text{et}\quad
g(i,\underline{x})>0
\quad\text{pour tout $0\leq i<z$}.
\]
\end{defn}

Informatiquement, $\mu g$ correspond à la recherche non bornée
\[
\texttt{i=0; while (true) \{ if (g(i,x)==0) return i; i++; \}}.
\]
Si aucune valeur convenable n'est trouvée, ou si l'évaluation de
$g(i,\underline{x})$ ne termine pas avant d'atteindre une éventuelle
solution, le calcul ne termine pas.

\begin{defn}\label{general-recursive-definition}
La classe $\mathbf{R}$ des \defin[générale récursive
(fonction)]{fonctions générales récursives}, ou simplement
\defin[récursive (fonction)]{fonctions récursives}, est la plus petite
classe de fonctions
$\mathbb{N}^k\dasharrow\mathbb{N}$, pour $k$ variable, qui~:
\begin{itemize}
\item contient les projections, les constantes et le successeur~;
\item est stable par composition~;
\item est stable par récursion primitive~;
\item est stable par l'opérateur $\mu$.
\end{itemize}
\end{defn}

Cette dernière clause autorise précisément les boucles dont le nombre
d'itérations n'est pas borné à l'avance.  Les fonctions ainsi obtenues
peuvent être partielles.

\thingy On numérote les fonctions générales récursives exactement
comme les fonctions primitives récursives, mais en ajoutant une clause
pour l'opérateur $\mu$.  Pour certains entiers $e$, on définit donc
\[
\varphi_e^{(k)}\colon\mathbb{N}^k\dasharrow\mathbb{N}.
\]
Les codes correspondant aux projections, constantes, successeur,
composition et récursion primitive sont les mêmes que précédemment.
On ajoute la clause suivante~: si
\[
e=\dbllangle5,k,c\dblrangle
\quad\text{et}\quad
g:=\varphi_c^{(k+1)},
\]
alors
\[
\varphi_e^{(k)}=\mu g.
\]
Les autres codes ne définissent aucune fonction.

Une fonction $f\colon\mathbb{N}^k\dasharrow\mathbb{N}$ est récursive
si et seulement s'il existe un indice $e$ tel que
\[
f=\varphi_e^{(k)}.
\]

\thingy Contrairement au cas primitif récursif, il existe cette fois
une fonction récursive universelle~:
\[
(e,\underline{x})\mapsto\varphi_e^{(k)}(\underline{x}).
\]
Autrement dit, il existe un indice $u_k$ tel que
\[
\varphi_{u_k}^{(k+1)}(e,\underline{x})
 =\varphi_e^{(k)}(\underline{x}).
\]
Informatiquement, cette fonction est un interpréteur du langage des
fonctions récursives écrit dans ce même langage.

Pour justifier cette propriété, on introduit des \defin{arbres de
calcul}.  Un arbre de calcul de
$\varphi_e^{(k)}(\underline{x})$ ayant pour résultat $y$ est un arbre
fini dont la racine affirme
\[
\varphi_e^{(k)}(\underline{x})=y
\]
et dont les sous-arbres justifient les calculs intermédiaires requis
par la clause de construction correspondant à $e$.

Pour une composition, les sous-arbres établissent les valeurs des
fonctions intérieures puis celle de la fonction extérieure.  Pour une
récursion primitive, ils établissent successivement les valeurs
précédentes.  Pour une minimisation $\mu g$, ils établissent les
valeurs positives
\[
g(0,\underline{x}),\ldots,g(y-1,\underline{x})
\]
et la valeur nulle $g(y,\underline{x})$.

On a alors
\[
\varphi_e^{(k)}(\underline{x})=y
\]
si et seulement s'il existe un arbre de calcul fini qui l'atteste.
Or vérifier qu'un entier code un arbre de calcul valable est une
opération primitive récursive.  Il suffit donc de parcourir les
entiers $n=0,1,2,\ldots$, de tester si $n$ code un arbre de calcul
valable pour $\varphi_e^{(k)}(\underline{x})$, et, dès que c'est le
cas, de renvoyer la valeur inscrite à la racine.  Cette recherche non
bornée utilise une seule application de l'opérateur $\mu$ et définit
la fonction universelle.

\begin{thm}[Forme normale de Kleene]
\label{normal-form-theorem}
Il existe un prédicat primitif récursif $T$ sur $\mathbb{N}^3$ et une
fonction primitive récursive $U\colon\mathbb{N}\to\mathbb{N}$ tels que
\[
\varphi_e^{(k)}(x_1,\ldots,x_k)
 =
U\bigl(\mu T(e,\dbllangle x_1,\ldots,x_k\dblrangle)\bigr).
\]
\end{thm}

Ici, $T(n,e,\dbllangle x_1,\ldots,x_k\dblrangle)$ exprime que $n$
code un arbre de calcul valable de
$\varphi_e^{(k)}(x_1,\ldots,x_k)$, tandis que $U$ extrait de cet arbre
son résultat.  Ce théorème affirme donc que toute fonction récursive
peut être calculée à l'aide d'une unique recherche non bornée, entourée
de calculs primitifs récursifs.

\thingy Le théorème s-m-n reste valable pour la numérotation
$\varphi_e^{(k)}$.

\begin{thm}[Théorème s-m-n]\label{smn-recursive}
Pour tous $m,n$, il existe une fonction primitive récursive
\[
s_{m,n}\colon\mathbb{N}^{m+1}\to\mathbb{N}
\]
telle que
\[
\varphi^{(n)}_{s_{m,n}(e,x_1,\ldots,x_m)}
       (y_1,\ldots,y_n)
 =
\varphi^{(m+n)}_e
       (x_1,\ldots,x_m,y_1,\ldots,y_n).
\]
\end{thm}

Le fait que les fonctions manipulées puissent maintenant être
partielles ne change pas la construction syntaxique du programme.
La fonction $s_{m,n}$ reste donc primitive récursive.

\thingy Le théorème de récursion de Kleene s'étend lui aussi sans
modification essentielle.

\begin{thm}[Kleene]\label{kleene-recursion-theorem}
Si
\[
h\colon\mathbb{N}^{k+1}\dasharrow\mathbb{N}
\]
est récursive, il existe un indice $e$ tel que
\[
(\forall\underline{x})\qquad
\varphi_e^{(k)}(\underline{x})=h(e,\underline{x}).
\]
De plus, on peut calculer primitivement récursivement un tel indice à
partir d'un indice de $h$ et de l'arité $k$.
\end{thm}

\begin{proof}
La preuve est la même que dans le cas primitif récursif.  Soit
$s:=s_{1,k}$.  La fonction
\[
(t,\underline{x})\mapsto h(s(t,t),\underline{x})
\]
est récursive~; écrivons-la sous la forme
$\varphi_c^{(k+1)}(t,\underline{x})$.  Alors
\[
\begin{aligned}
\varphi_{s(c,c)}^{(k)}(\underline{x})
 &=\varphi_c^{(k+1)}(c,\underline{x})\\
 &=h(s(c,c),\underline{x}).
\end{aligned}
\]
L'indice $e:=s(c,c)$ convient.
\end{proof}

Une reformulation utile fait intervenir l'universalité.

\begin{thm}[Kleene--Rogers]\label{kleene-rogers-fixed-point}
Si $F\colon\mathbb{N}\to\mathbb{N}$ est récursive et si
$k\in\mathbb{N}$, il existe un indice $e$ tel que
\[
\varphi_e^{(k)}=\varphi_{F(e)}^{(k)}.
\]
\end{thm}

Ainsi, pour toute transformation calculable $F$ appliquée au code
source d'un programme, il existe un programme qui calcule la même
fonction que son transformé par $F$.

\thingy Malgré leur nom, les fonctions générales récursives ne sont pas
définies directement par des appels récursifs.  Le langage de la
définition~\ref{general-recursive-definition} ne contient que des
opérations élémentaires, la composition, la récursion primitive et la
recherche non bornée.  Le théorème de récursion permet néanmoins
d'implémenter les appels récursifs~: on donne à la fonction en cours de
définition accès à son propre indice, puis chaque appel récursif est
remplacé par un appel à la fonction universelle sur cet indice.

\subsection{Le problème de l'arrêt}
\label{subsection-halting-problem-recursive-functions}

\thingy Étant donné un programme d'indice $e$ et une entrée $x$, peut-on
déterminer algorithmiquement si le calcul
$\varphi_e^{(1)}(x)$ termine~?  La réponse est négative.

\begin{thm}[Turing]\label{undecidability-halting-problem}
Il n'existe pas de fonction récursive totale
$h\colon\mathbb{N}^2\to\mathbb{N}$ telle que
\[
h(e,x)=
\begin{cases}
1&\text{si $\varphi_e^{(1)}(x)\downarrow$,}\\
0&\text{si $\varphi_e^{(1)}(x)\uparrow$.}
\end{cases}
\]
\end{thm}

\begin{proof}
Supposons par l'absurde qu'une telle fonction $h$ existe.  Définissons
la fonction partielle
\[
v(e,x):=
\begin{cases}
42&\text{si $h(e,x)=0$,}\\
\uparrow&\text{si $h(e,x)=1$.}
\end{cases}
\]
Cette fonction est récursive~: on calcule $h(e,x)$, puis on renvoie
$42$ dans le premier cas et on effectue une recherche infinie dans le
second.

Par le théorème de récursion de Kleene, il existe un indice $e$ tel que
\[
\varphi_e^{(1)}(x)=v(e,x).
\]
Si $\varphi_e^{(1)}(x)\downarrow$, alors $h(e,x)=1$, donc
$v(e,x)\uparrow$, ce qui contredit l'égalité précédente.  Si
$\varphi_e^{(1)}(x)\uparrow$, alors $h(e,x)=0$, donc
$v(e,x)\downarrow$, ce qui donne encore une contradiction.
\end{proof}

Une version plus directement diagonale consiste à supposer que $h$
décide l'arrêt et à définir
\[
v(e):=
\begin{cases}
42&\text{si $h(e,e)=0$,}\\
\uparrow&\text{si $h(e,e)=1$.}
\end{cases}
\]
Si $v=\varphi_c^{(1)}$, alors l'étude de $\varphi_c^{(1)}(c)$ conduit
immédiatement à une contradiction~: elle termine si et seulement si
elle ne termine pas.

\thingy Les fonctions primitives récursives sont toutes totales, de
sorte que leur problème de l'arrêt est trivial.  Elles ne possèdent en
revanche pas de fonction universelle primitive récursive.  Les
fonctions générales récursives peuvent être partielles et disposent
d'une fonction universelle, mais leur problème de l'arrêt est
indécidable.  L'universalité est précisément ce qui permet de
construire l'argument diagonal.

\subsection{Machines de Turing}
\label{subsection-turing-machines}

\thingy Une \defin[Turing (machine de)]{machine de Turing} est une
modélisation d'un ordinateur extrêmement simple.  Elle possède un état
interne choisi dans un ensemble fini et une mémoire illimitée sous la
forme d'une bande linéaire divisée en cellules.  Chaque cellule contient
un symbole pris dans un alphabet fini.  Une tête de lecture et
d'écriture se trouve au-dessus d'une unique cellule à chaque instant.

Le programme de la machine indique, en fonction de l'état courant et
du symbole lu par la tête~:
\begin{itemize}
\item le nouvel état~;
\item le symbole à écrire dans la cellule courante~;
\item la direction dans laquelle déplacer la tête, à gauche ou à
  droite.
\end{itemize}
La machine poursuit son exécution jusqu'à atteindre un état spécial
d'arrêt.

\begin{defn}\label{definition-turing-machine}
Une machine de Turing déterministe à une bande, deux symboles et
$m\geq2$ états est la donnée~:
\begin{itemize}
\item d'un ensemble fini $Q$ de cardinal $m$, identifié à
  $\{0,\ldots,m-1\}$~;
\item d'un ensemble de symboles de bande
  $\Sigma=\{0,1\}$~;
\item d'une fonction
\[
\delta\colon
(Q\setminus\{0\})\times\Sigma
 \longrightarrow
Q\times\Sigma\times\{\texttt{L},\texttt{R}\},
\]
appelée le programme de la machine.
\end{itemize}
L'état $0$ est l'état d'arrêt.
\end{defn}

Une \defin{configuration} d'une telle machine est un triplet
$(q,\beta,i)$, où~:
\begin{itemize}
\item $q\in Q$ est l'état courant~;
\item $\beta\colon\mathbb{Z}\to\Sigma$ est la bande, supposée ne prendre
  qu'un nombre fini de valeurs différentes de $0$~;
\item $i\in\mathbb{Z}$ est la position de la tête.
\end{itemize}

Si $q\neq0$ et si
\[
(q',y,d)=\delta(q,\beta(i)),
\]
la configuration suivante est $(q',\beta',i')$, où
\[
i'=
\begin{cases}
i-1&\text{si $d=\texttt{L}$,}\\
i+1&\text{si $d=\texttt{R}$,}
\end{cases}
\]
et
\[
\beta'(j)=
\begin{cases}
y&\text{si $j=i$,}\\
\beta(j)&\text{si $j\neq i$.}
\end{cases}
\]

\thingy À partir d'une configuration initiale $C$, la
\defin[exécution (trace d')]{trace d'exécution} est la suite finie ou
infinie
\[
C^{(0)},C^{(1)},C^{(2)},\ldots
\]
où $C^{(0)}=C$ et où $C^{(n+1)}$ est la configuration suivant
$C^{(n)}$, tant que l'état courant n'est pas $0$.  Si une configuration
$C^{(n)}$ d'état $0$ est atteinte, la machine s'arrête après $n$
étapes et $C^{(n)}$ est sa configuration finale.

\thingy Une machine de Turing peut être codée par un entier, puisque
son programme est une table finie.  Une configuration peut elle aussi
être codée par un entier, puisque la bande ne contient qu'un nombre
fini de symboles non nuls.

La fonction
\[
(M,C)\mapsto C'
\]
associant à une machine et à une configuration la configuration
suivante est primitive récursive.  Par récursion primitive, la
fonction
\[
(n,M,C)\mapsto C^{(n)}
\]
associant la configuration atteinte après $n$ étapes l'est également.

En revanche, la fonction qui associe à $(M,C)$ sa configuration finale
si la machine s'arrête, et qui est non définie sinon, est générale
récursive.  En effet, on peut rechercher sans borne le premier entier
$n$ tel que $C^{(n)}$ soit une configuration d'arrêt.

\thingy Pour définir précisément les fonctions calculées par les
machines de Turing, fixons une convention d'entrée et de sortie.  Une
fonction
\[
f\colon\mathbb{N}^k\dasharrow\mathbb{N}
\]
est dite \defin[calculable par machine de Turing]{calculable par
machine de Turing} lorsqu'il existe une machine qui, sur les entrées
$x_1,\ldots,x_k$, part d'une configuration initiale dans laquelle~:
\begin{itemize}
\item l'état vaut $1$ et la tête se trouve en position $0$~;
\item la partie négative du ruban est laissée intacte~;
\item la partie positive contient le mot
\[
0\,1^{x_1}\,0\,1^{x_2}\,0\cdots0\,1^{x_k}\,0,
\]
suivi d'une infinité de $0$~;
\end{itemize}
et telle que~:
\begin{itemize}
\item si $f(x_1,\ldots,x_k)\uparrow$, la machine ne s'arrête pas~;
\item si $f(x_1,\ldots,x_k)\downarrow=y$, elle s'arrête avec la tête en
  position $0$, sans modifier la partie négative du ruban, et avec
  \[
  0\,1^y\,0
  \]
  sur la partie positive du ruban, suivi d'une infinité de $0$.
\end{itemize}

Ces conventions sont arbitraires mais sans importance pour la classe
de fonctions obtenue.

\begin{thm}\label{recursive-functions-turing-machines}
Une fonction
$f\colon\mathbb{N}^k\dasharrow\mathbb{N}$ est générale récursive si et
seulement si elle est calculable par une machine de Turing.
\end{thm}

\begin{proof}
Pour simuler les machines de Turing par les fonctions récursives, on
code les programmes et les configurations par des entiers.  La
fonction donnant la configuration après $n$ étapes est primitive
récursive.  Une recherche non bornée permet de trouver le premier
$n$ pour lequel la machine est à l'arrêt, puis une fonction primitive
récursive décode la sortie.

Réciproquement, on montre par induction sur la
définition~\ref{general-recursive-definition} que toute fonction
générale récursive est calculable par machine de Turing.  Il faut
programmer les projections, les constantes et le successeur, puis
montrer que les fonctions calculables par machine de Turing sont
stables par composition, récursion primitive et minimisation.  Ces
constructions reviennent à sauvegarder les entrées et les résultats
intermédiaires sur la bande, puis à exécuter successivement les
machines correspondant aux sous-calculs.
\end{proof}

\review{Les diapositives qualifient la seconde direction de
«~fastidieuse mais pas difficile~» et n'en donnent pas tous les détails.
La preuve ci-dessus reste donc une esquisse à développer ou à renvoyer
à un exercice.}

Cette équivalence est constructive~: on peut calculer un code de
machine de Turing à partir d'un indice de fonction récursive, et
réciproquement.

\thingy De nombreuses variantes de la définition sont possibles~:
machines à plusieurs bandes, plusieurs têtes, davantage de symboles ou
transitions non déterministes.  Du point de vue de la calculabilité,
ces variantes ne rendent pas le modèle plus puissant, sauf cas
dégénérés.  Elles peuvent cependant modifier considérablement la
complexité des simulations.

Une machine de Turing universelle peut recevoir sur sa bande le code
d'une autre machine $M$ et d'une configuration initiale $C$, puis
simuler l'exécution de $M$ sur $C$.  Elle s'arrête si et seulement si
$M$ s'arrête.  Cette propriété est la traduction, pour les machines de
Turing, de l'existence d'une fonction récursive universelle.

\begin{thm}\label{halting-problem-turing-machines}
Il n'existe pas d'algorithme qui, prenant en entrée une machine de
Turing $M$ et une configuration $C$, décide si $M$ s'arrête à partir
de $C$.
\end{thm}

\begin{proof}
Si un tel algorithme existait, on pourrait coder $M$ et $C$ par des
entiers et décider le problème de l'arrêt pour les fonctions
récursives, en contradiction avec le
théorème~\ref{undecidability-halting-problem}.
\end{proof}

Le problème reste indécidable si l'on demande seulement si une machine
s'arrête à partir d'une bande vierge.  En effet, étant donnés $M$ et
$C$, le théorème s-m-n permet de construire une machine $M^*$ qui
prépare d'abord la configuration $C$ à partir de la bande vierge, puis
simule $M$.  La machine $M^*$ s'arrête sur bande vierge si et seulement
si $M$ s'arrête sur $C$.

\thingy La \defin[castor affairé]{fonction du castor affairé} associe à
$m$ le nombre maximal $B(m)$ d'étapes effectuées, à partir d'une bande
vierge, par une machine de Turing ayant au plus $m$ états et qui finit
par s'arrêter.  Ce maximum existe, puisqu'il n'existe qu'un nombre fini
de machines ayant au plus $m$ états.

\begin{prop}\label{busy-beaver-not-computable}
La fonction $B$ n'est pas calculable.
\end{prop}

\begin{proof}
Supposons qu'il existe une fonction calculable
$f\colon\mathbb{N}\to\mathbb{N}$ telle que
\[
B(m)\leq f(m)
\]
pour tout $m$.  Étant donnée une machine $M$ à $m$ états, on pourrait
l'exécuter pendant $f(m)$ étapes.  Si elle s'arrête pendant ce temps,
on répondrait qu'elle termine.  Sinon, la définition de $B(m)$ assure
qu'elle ne terminera jamais.  On aurait ainsi un algorithme décidant
le problème de l'arrêt sur bande vierge, ce qui est impossible.
\end{proof}

La fonction du castor affairé croît donc plus vite que toute fonction
calculable, au sens suivant~: pour toute fonction calculable totale
$f\colon\mathbb{N}\to\mathbb{N}$, il existe $m_0$ tel que
\[
B(m)>f(m)
\]
pour tout $m\geq m_0$.

\review{Les diapositives esquissent une preuve de cette domination à
l'aide d'une famille de machines dont le nombre d'états est contrôlé.
La formulation uniforme et les détails de codage devront être
vérifiés.}

\subsection{Décidabilité et semi-décidabilité}
\label{subsection-decidability-semidecidability}

\thingy D'après les équivalences précédentes, une fonction partielle
est calculable si et seulement si elle est générale récursive, si et
seulement si elle est calculable par machine de Turing.  Nous
emploierons librement ces différentes descriptions.

On rencontre aussi les synonymes suivants~:
\begin{itemize}
\item une fonction partielle calculable est parfois dite
  «~semi-calculable~» ou «~générale récursive~»~;
\item un ensemble décidable est parfois dit «~calculable~» ou
  «~récursif~»~;
\item un ensemble semi-décidable est parfois dit
  «~semi-calculable~», «~calculablement énumérable~» ou
  «~récursivement énumérable~».
\end{itemize}

\begin{prop}\label{decidable-iff-semidecidable-and-complement}
Une partie $A\subseteq\mathbb{N}^k$ est décidable si et seulement si
$A$ et son complémentaire
\[
\complement A:=\mathbb{N}^k\setminus A
\]
sont tous deux semi-décidables.
\end{prop}

\begin{proof}
Si $A$ est décidable, son complémentaire est décidable en échangeant
les réponses $0$ et $1$.  De plus, tout ensemble décidable est
semi-décidable~: il suffit d'exécuter son algorithme de décision et de
boucler indéfiniment lorsque celui-ci répond «~non~».

Réciproquement, supposons que $A$ et $\complement A$ soient
semi-décidables.  Soient $e_1$ et $e_2$ des indices tels que
\[
\varphi_{e_1}(\underline{x})\downarrow
 \quad\Longleftrightarrow\quad
\underline{x}\in A
\]
et
\[
\varphi_{e_2}(\underline{x})\downarrow
 \quad\Longleftrightarrow\quad
\underline{x}\notin A.
\]
On exécute les deux calculs en parallèle, en alternant leurs étapes.
Exactement l'un des deux finit par terminer.  S'il s'agit du premier,
on répond «~oui~»~; s'il s'agit du second, on répond «~non~».
L'algorithme ainsi obtenu termine sur toute entrée et décide $A$.
\end{proof}

Le «~lancement en parallèle~» ne nécessite aucun véritable parallélisme.
On peut simuler successivement une étape du premier programme, puis une
étape du second, et recommencer jusqu'à ce que l'un d'eux s'arrête.

\thingy Le \defin[problème de l'arrêt]{problème de l'arrêt} est
l'ensemble
\[
\mathscr{H}
 :=\{(e,x)\in\mathbb{N}^2:
           \varphi_e^{(1)}(x)\downarrow\}.
\]
Il n'est pas décidable, d'après le
théorème~\ref{undecidability-halting-problem}.  Il est en revanche
semi-décidable~: sur l'entrée $(e,x)$, il suffit d'exécuter
$\varphi_e^{(1)}(x)$ et de répondre «~oui~» si le calcul termine.

La proposition~\ref{decidable-iff-semidecidable-and-complement} montre
donc que son complémentaire $\complement\mathscr{H}$ n'est pas
semi-décidable.  Il est possible de reconnaître en temps fini un
calcul qui termine, en observant effectivement sa terminaison, mais il
n'existe aucun procédé analogue reconnaissant tous les calculs qui ne
termineront jamais.

\thingy Si $A\subseteq\mathbb{N}$ est décidable et si
$f\colon\mathbb{N}\to\mathbb{N}$ est calculable totale, alors l'image
\[
f(A):=\{f(i):i\in A\}
\]
est semi-décidable.

En effet, pour semi-décider si $m\in f(A)$, on parcourt
$i=0,1,2,\ldots$.  Pour chaque $i$, on décide si $i\in A$ et, si c'est
le cas, on calcule $f(i)$.  Si $f(i)=m$, on termine en répondant
«~oui~»~; sinon, on poursuit la recherche.

Réciproquement, tout ensemble semi-décidable $B\subseteq\mathbb{N}$
peut être obtenu comme image d'un ensemble décidable par une fonction
totale calculable.

Une autre formulation, lorsque $B$ est non vide, est la suivante~:
$B$ est semi-décidable si et seulement s'il existe une fonction totale
calculable $f\colon\mathbb{N}\to\mathbb{N}$ telle que
\[
f(\mathbb{N})=B.
\]
C'est ce qui justifie la terminologie «~calculablement énumérable~».

\thingy Les ensembles décidables sont stables par réunions finies,
intersections finies et complémentaire.  Ils ne sont cependant pas
stables par projection.  Le problème de l'arrêt en donne déjà un
exemple, puisqu'il peut être présenté comme la projection d'un ensemble
décidable exprimant qu'un calcul s'arrête en au plus $n$ étapes.

Les ensembles semi-décidables sont stables par réunions finies,
intersections finies et projections.  Ils ne sont pas stables par
complémentaire, puisque $\mathscr{H}$ est semi-décidable tandis que
$\complement\mathscr{H}$ ne l'est pas.

\subsection{Le théorème de Rice}
\label{subsection-rice-theorem}

\thingy Soit $\mathbf{R}^{(1)}$ l'ensemble des fonctions partielles
calculables $\mathbb{N}\dasharrow\mathbb{N}$, et considérons
l'application
\[
\Phi\colon e\mapsto\varphi_e^{(1)}.
\]
Cette application associe à un programme, son intention, la fonction
qu'il calcule, son extension.

\begin{thm}[Rice]\label{rice-theorem}
Soit $F\subseteq\mathbf{R}^{(1)}$ un ensemble de fonctions partielles
calculables.  Si
\[
\Phi^{-1}(F)
 =\{e\in\mathbb{N}:\varphi_e^{(1)}\in F\}
\]
est décidable, alors
\[
F=\varnothing
\qquad\text{ou}\qquad
F=\mathbf{R}^{(1)}.
\]
\end{thm}

Autrement dit, aucune propriété non triviale de la fonction calculée
par un programme n'est décidable en examinant le programme.  Par
exemple, aucun des ensembles suivants n'est décidable~:
\[
\{e:\varphi_e(0)\downarrow\},
\]
\[
\{e:\varphi_e\text{ est totale}\},
\]
ou encore
\[
\{e:(\forall n)\,
       (\varphi_e(n)\downarrow\limp\varphi_e(n)=0)\}.
\]

\begin{proof}
Supposons par l'absurde que $\Phi^{-1}(F)$ soit décidable, avec
$F\neq\varnothing$ et $F\neq\mathbf{R}^{(1)}$.  Choisissons
$f\in F$ et $g\notin F$.  Définissons
\[
h(e,x):=
\begin{cases}
f(x)&\text{si $e\notin\Phi^{-1}(F)$,}\\
g(x)&\text{si $e\in\Phi^{-1}(F)$.}
\end{cases}
\]
La fonction $h$ est calculable, puisque l'appartenance de $e$ à
$\Phi^{-1}(F)$ est supposée décidable.

Par le théorème de récursion de Kleene, il existe un indice $e$ tel
que
\[
\varphi_e^{(1)}(x)=h(e,x).
\]
Si $e\in\Phi^{-1}(F)$, alors $h(e,x)=g(x)$ pour tout $x$, donc
$\varphi_e^{(1)}=g\notin F$, contradiction.  Si
$e\notin\Phi^{-1}(F)$, alors $\varphi_e^{(1)}=f\in F$, contradiction
également.
\end{proof}

\thingy Le théorème de Rice peut aussi être démontré par réduction au
problème de l'arrêt.  Supposons $F$ non trivial et, quitte à remplacer
$F$ par son complémentaire, supposons que la fonction nulle part
définie n'appartienne pas à $F$.  Choisissons $f\in F$ et un indice
$a$ de $f$.

Pour $(e,x)\in\mathbb{N}^2$, construisons un programme d'indice
$b(e,x)$ qui commence par simuler $\varphi_e(x)$.  Si cette simulation
termine, il calcule ensuite $f$ sur son argument.  Ainsi,
\[
\varphi_{b(e,x)}=
\begin{cases}
f&\text{si $\varphi_e(x)\downarrow$,}\\
{\uparrow}&\text{si $\varphi_e(x)\uparrow$,}
\end{cases}
\]
où ${\uparrow}$ désigne la fonction nulle part définie.  On obtient
donc
\[
b(e,x)\in\Phi^{-1}(F)
\quad\Longleftrightarrow\quad
(e,x)\in\mathscr{H}.
\]
Si $\Phi^{-1}(F)$ était décidable, le problème de l'arrêt le serait
aussi, ce qui est impossible.

\subsection{Réductions}
\label{subsection-reductions}

\thingy Pour montrer qu'un problème $A$ est indécidable, on cherche
souvent à transformer toute instance d'un problème déjà connu comme
indécidable, par exemple $\mathscr{H}$, en une instance de $A$.  Cette
transformation est appelée une \defin{réduction}.

\begin{defn}\label{definition-many-one-reduction}
Si $A,B\subseteq\mathbb{N}$, on écrit
\[
A\mathrel{\leq_{\mathrm m}}B
\]
lorsqu'il existe une fonction totale calculable
$\rho\colon\mathbb{N}\to\mathbb{N}$ telle que
\[
\rho(m)\in B
\quad\Longleftrightarrow\quad
m\in A.
\]
On parle de \defin[réduction many-to-one]{réduction
\textit{many-to-one}}.
\end{defn}

Ainsi, pour décider si $m\in A$, il suffirait de calculer
$\rho(m)$ puis de décider si $\rho(m)\in B$.  Par conséquent, si
$A\leq_{\mathrm m}B$ et si $B$ est décidable, alors $A$ est décidable.
La même implication vaut pour la semi-décidabilité.  En particulier,
si
\[
\mathscr{H}\leq_{\mathrm m}D,
\]
alors $D$ n'est pas décidable.

La relation $\leq_{\mathrm m}$ est réflexive et transitive.  La relation
\[
A\equiv_{\mathrm m}B
\quad\Longleftrightarrow\quad
A\leq_{\mathrm m}B
\ \text{et}\
B\leq_{\mathrm m}A
\]
est donc une relation d'équivalence.

\thingy Une notion plus générale autorise un algorithme à interroger
plusieurs fois un oracle décidant $B$ au cours de son exécution.

Informellement, on écrit
\[
A\leq_{\mathrm T}B
\]
lorsqu'il existe un algorithme qui décide $A$ et qui peut, au cours de
son calcul, demander si un entier $n$ appartient à $B$.  Cette notion
est appelée \defin[Turing (réduction de)]{réduction de Turing}.

La réduction de Turing est plus faible que la réduction many-to-one~:
une réduction many-to-one ne pose qu'une question à l'oracle, à la fin
du calcul.  Une réduction de Turing peut en poser un nombre arbitraire,
mais fini, et choisir chaque nouvelle question en fonction des réponses
précédentes.

\review{Les diapositives proposent plusieurs formalisations possibles
des réductions de Turing.  Leur définition entièrement formelle devra
être harmonisée avec la présentation retenue pour les oracles dans le
document final.}

Si $A\leq_{\mathrm T}B$ et si $B$ est décidable, alors $A$ est
décidable.  Il s'ensuit que si $\mathscr{H}\leq_{\mathrm T}D$, alors
$D$ est indécidable.

\subsection{Le $\lambda$-calcul non typé}
\label{subsection-untyped-lambda-calculus}

\thingy Le \defin[$\lambda$-calcul]{$\lambda$-calcul non typé}
manipule des expressions telles que
\[
\lambda x.\lambda y.\lambda z.((xz)(yz)),
\qquad
\lambda f.\lambda x.f(f(f(f(fx)))),
\]
ou
\[
(\lambda x.(xx))(\lambda x.(xx)).
\]
Ces expressions, appelées \defin[$\lambda$-terme]{$\lambda$-termes},
représentent intuitivement des fonctions prenant d'autres fonctions en
entrée et renvoyant des fonctions.

Les deux constructions fondamentales sont~:
\begin{itemize}
\item l'\defin[application ($\lambda$-calcul)]{application} $(PQ)$,
  qui consiste à appliquer la fonction représentée par $P$ à
  l'argument représenté par $Q$~;
\item l'\defin[abstraction ($\lambda$-calcul)]{abstraction}
  $\lambda v.E$, qui représente la fonction associant à un argument
  la valeur obtenue en remplaçant $v$ par cet argument dans $E$.
\end{itemize}

\begin{defn}\label{definition-lambda-terms}
Un terme du $\lambda$-calcul est défini inductivement de la manière
suivante~:
\begin{itemize}
\item toute variable est un terme~;
\item si $P$ et $Q$ sont des termes, alors $(PQ)$ est un terme~;
\item si $v$ est une variable et $E$ un terme, alors $\lambda v.E$ est
  un terme.
\end{itemize}
Dans $\lambda v.E$, les occurrences de $v$ dans $E$ auxquelles ce
$\lambda$ s'applique sont dites \defin[liée
(variable)]{liées}.  Les autres occurrences de variables sont dites
\defin[libre (variable)]{libres}.  Un terme sans variable libre est dit
\defin[clos ($\lambda$-terme)]{clos}.
\end{defn}

Par convention, l'application est parenthésée vers la gauche~:
\[
xyz:=((xy)z).
\]
On abrège
\[
\lambda u.\lambda v.E
\]
en $\lambda uv.E$.  Enfin, l'abstraction est moins prioritaire que
l'application~:
\[
\lambda x.xy
\]
signifie $\lambda x.(xy)$ et non $(\lambda x.x)y$.

\thingy Le nom d'une variable liée ne doit pas avoir d'importance.  On
appelle \defin[$\alpha$-conversion]{$\alpha$-conversion} le renommage
cohérent d'une variable liée.  Ainsi,
\[
\lambda x.x\equiv\lambda y.y.
\]
Il faut cependant éviter de capturer une variable libre.  Par exemple,
\[
\lambda y.xy\not\equiv\lambda x.xx.
\]
Dans une expression comportant plusieurs abstractions portant le même
nom, une occurrence est liée par le $\lambda$ le plus intérieur
approprié.

Une façon de supprimer complètement les noms de variables liées
consiste à utiliser les indices de De Bruijn~: chaque occurrence liée
est remplacée par la distance qui la sépare de l'abstraction qui la
lie.  Deux termes sont alors $\alpha$-équivalents si et seulement si
leurs écritures à l'aide de ces indices sont identiques.

\begin{defn}\label{definition-beta-reduction}
Un \defin{redex} est un terme de la forme
\[
(\lambda v.E)T.
\]
Son \defin[réduit ($\lambda$-calcul)]{réduit} est le terme
\[
E[v\backslash T]
\]
obtenu en remplaçant dans $E$ les occurrences libres de $v$ par $T$,
en renommant au besoin des variables liées pour éviter toute capture.

Le remplacement d'un redex par son réduit est appelé une
\defin[$\beta$-réduction]{$\beta$-réduction}.
\end{defn}

Par exemple,
\[
(\lambda x.xx)y\longrightarrow yy,
\]
tandis que
\[
(\lambda x.xx)(\lambda x.xx)
 \longrightarrow
(\lambda x.xx)(\lambda x.xx),
\]
de sorte que ce dernier terme se réduit indéfiniment.

On note $T\longrightarrow T'$ lorsque $T'$ est obtenu par une seule
$\beta$-réduction, et
\[
T\twoheadrightarrow T'
\]
lorsque $T'$ est obtenu par une suite finie, éventuellement vide, de
$\beta$-réductions.

Un terme ne contenant aucun redex est dit en
\defin[forme normale ($\lambda$-calcul)]{forme normale}.  Il est
\defin[faiblement normalisable]{$\beta$-faiblement normalisable}
lorsqu'au moins une suite de réductions mène à une forme normale, et
\defin[fortement normalisable]{$\beta$-fortement normalisable} lorsque
toute suite de réductions termine.

\begin{thm}[Church--Rosser]\label{church-rosser-theorem}
Si
\[
T\twoheadrightarrow T'_1
\quad\text{et}\quad
T\twoheadrightarrow T'_2,
\]
alors il existe $T''$ tel que
\[
T'_1\twoheadrightarrow T''
\quad\text{et}\quad
T'_2\twoheadrightarrow T''.
\]
\end{thm}

En particulier, si $T'_1$ et $T'_2$ sont en forme normale, ils sont
$\alpha$-équivalents.  La forme normale, lorsqu'elle existe, est donc
unique à $\alpha$-conversion près.

\thingy Pour éviter d'avoir à démontrer ici le théorème de
Church--Rosser, on choisit une stratégie déterministe de réduction.  Le
\defin[redex extérieur gauche]{redex extérieur gauche} est le redex
dont le symbole $\lambda$ est situé le plus à gauche dans l'écriture du
terme.  On note
\[
T\longrightarrow_{\mathsf{lft}}T'
\]
la réduction de ce redex, et
\[
T\twoheadrightarrow_{\mathsf{lft}}T'
\]
une suite de telles réductions.

\begin{thm}[Curry et al.]\label{leftmost-normalization}
Si un terme $T$ possède une forme normale $T'$, alors
\[
T\twoheadrightarrow_{\mathsf{lft}}T'.
\]
\end{thm}

La stratégie extérieure gauche est donc déterministe et trouve une
forme normale chaque fois qu'une forme normale peut être atteinte.

\subsection{Entiers de Church et fonctions représentables}
\label{subsection-church-numerals}

\thingy Pour chaque entier $n\in\mathbb{N}$, on définit l'\defin[Church
(entier de)]{entier de Church}
\[
\overline{n}:=\lambda fx.f^{\circ n}(x).
\]
Ainsi,
\[
\begin{aligned}
\overline{0}&:=\lambda fx.x,\\
\overline{1}&:=\lambda fx.fx,\\
\overline{2}&:=\lambda fx.f(fx),\\
\overline{3}&:=\lambda fx.f(f(fx)).
\end{aligned}
\]
Intuitivement, $\overline{n}$ reçoit une fonction $f$ et un argument
$x$, puis applique $f$ exactement $n$ fois à $x$.

Posons
\[
A:=\lambda mfx.f(mfx).
\]
Alors
\[
A\overline{n}
 \twoheadrightarrow_{\mathsf{lft}}
\overline{n+1},
\]
de sorte que $A$ représente la fonction successeur.

\begin{defn}\label{definition-lambda-representable}
Une fonction
\[
f\colon\mathbb{N}^k\dasharrow\mathbb{N}
\]
est \defin[représentable dans le $\lambda$-calcul]{représentable par un
$\lambda$-terme} lorsqu'il existe un terme clos $t$ tel que, pour tous
$x_1,\ldots,x_k\in\mathbb{N}$~:
\begin{itemize}
\item si $f(x_1,\ldots,x_k)\downarrow=y$, alors
\[
t\overline{x_1}\cdots\overline{x_k}
 \twoheadrightarrow_{\mathsf{lft}}\overline{y};
\]
\item si $f(x_1,\ldots,x_k)\uparrow$, alors la réduction extérieure
  gauche de
  $t\overline{x_1}\cdots\overline{x_k}$ ne termine pas.
\end{itemize}
\end{defn}

À titre d'exemple,
\[
\lambda mfx.f(mfx)
\]
représente le successeur, tandis que
\[
\lambda mnfx.nf(mfx)
\]
représente l'addition, et
\[
\lambda mnf.n(mf)
\]
représente la multiplication.

On peut également représenter la fonction
\[
(m,n,p)\mapsto
\begin{cases}
m&\text{si $p=0$,}\\
n&\text{si $p\geq1$.}
\end{cases}
\]

\thingy Les projections, les constantes, le successeur et la
composition sont aisément représentables dans le $\lambda$-calcul.  La
représentation de la récursion primitive exige de coder des couples
d'entiers.

On peut poser
\[
\begin{aligned}
\overline{m,n}
 &:=\lambda fgx.f^{\circ m}(g^{\circ n}(x)),\\
\Pi&:=\lambda mnfgx.(mf)(ngx),\\
\pi_1&:=\lambda pfx.pf(\lambda z.z)x,\\
\pi_2&:=\lambda pgx.p(\lambda z.z)gx.
\end{aligned}
\]
Alors
\[
\Pi\overline{m}\,\overline{n}
 \twoheadrightarrow_{\mathsf{lft}}\overline{m,n},
\]
et
\[
\pi_1\overline{m,n}
 \twoheadrightarrow_{\mathsf{lft}}\overline{m},
\qquad
\pi_2\overline{m,n}
 \twoheadrightarrow_{\mathsf{lft}}\overline{n}.
\]

À l'aide de ces couples, on peut faire évoluer simultanément la valeur
d'une récursion primitive et son compteur.  Ceci permet de représenter
toute fonction primitive récursive par un $\lambda$-terme.

\review{La formule explicite donnée dans les diapositives pour
représenter la récursion primitive est longue et dépend des conventions
précises de codage des couples.  Elle devra être vérifiée caractère par
caractère avant d'être intégrée dans une version définitive.}

\thingy Pour représenter toutes les fonctions générales récursives, il
reste à implémenter des appels récursifs, ou de façon équivalente
l'opérateur $\mu$.  On utilise pour cela le
\defin[Curry (combinateur $\mathsf{Y}$ de)]{combinateur de point fixe
de Curry}
\[
\mathsf{Y}
 :=\lambda f.((\lambda x.f(xx))(\lambda x.f(xx))).
\]
En effet,
\[
\begin{aligned}
\mathsf{Y}f
 &\longrightarrow
 f((\lambda x.f(xx))(\lambda x.f(xx)))\\
 &=f(\mathsf{Y}f),
\end{aligned}
\]
à réduction près.  Le terme $\mathsf{Y}f$ se comporte donc comme un
point fixe de $f$.

Si l'on veut définir une fonction récursive par
\[
\texttt{let rec }h=E,
\]
où $E$ contient des occurrences de $h$, on remplace cette définition
par un terme construit à partir de
\[
\mathsf{Y}(\lambda h.E).
\]
Le combinateur $\mathsf{Y}$ fournit ainsi au corps $E$ une
représentation de la fonction en cours de définition.

Le fonctionnement de $\mathsf{Y}$ dépend toutefois de la stratégie de
réduction.  La variante
\[
\mathsf{Z}
 :=\lambda f.
 ((\lambda x.f(\lambda v.xxv))
  (\lambda x.f(\lambda v.xxv)))
\]
maintient le point fixe sous une abstraction et convient mieux aux
langages qui évaluent les arguments avant l'appel d'une fonction.

\thingy L'opérateur $\mu$ peut être représenté en programmant la
recherche récursive
\[
h(z,x_1,\ldots,x_k)=
\begin{cases}
z&\text{si $g(z,x_1,\ldots,x_k)=0$,}\\
h(z+1,x_1,\ldots,x_k)
 &\text{sinon.}
\end{cases}
\]
Alors
\[
\mu g(x_1,\ldots,x_k)=h(0,x_1,\ldots,x_k).
\]
Le combinateur de point fixe permet de représenter $h$, et donc
$\mu g$, dans le $\lambda$-calcul.

Réciproquement, un $\lambda$-terme peut être codé par un entier.  Le
test déterminant si un terme est en forme normale, ainsi que la
fonction calculant sa réduction extérieure gauche, sont primitifs
récursifs.  Par récursion primitive, on peut calculer le terme obtenu
après $n$ réductions.  Une recherche non bornée permet alors de trouver
la première forme normale, lorsqu'elle existe.  Les fonctions
récursives peuvent donc simuler le $\lambda$-calcul.

\begin{thm}\label{recursive-functions-lambda-calculus}
Une fonction
$f\colon\mathbb{N}^k\dasharrow\mathbb{N}$ est générale récursive si et
seulement si elle est représentable par un terme du $\lambda$-calcul
non typé.
\end{thm}

\begin{proof}
Toute fonction générale récursive est représentable~: les fonctions
initiales et la composition sont directement représentables, la
récursion primitive l'est à l'aide du codage des couples, et la
minimisation l'est à l'aide d'un combinateur de point fixe.

Réciproquement, on code les termes par des entiers.  Les étapes de
réduction extérieure gauche sont primitives récursives et une
recherche non bornée trouve la première forme normale.  La lecture
d'un entier de Church dans cette forme normale est elle aussi
effective.  La fonction représentée est donc générale récursive.
\end{proof}

Cette équivalence est constructive~: on peut traduire
algorithmiquement un indice de fonction récursive en un
$\lambda$-terme qui la représente, et réciproquement.

\subsection{Conclusion}\label{subsection-computability-conclusion}

\thingy Nous avons obtenu trois descriptions équivalentes des fonctions
calculables.  Pour une fonction partielle
\[
f\colon\mathbb{N}^k\dasharrow\mathbb{N},
\]
les propriétés suivantes sont équivalentes~:
\begin{enumerate}
\item $f$ est générale récursive~;
\item $f$ est représentable dans le $\lambda$-calcul non typé~;
\item $f$ est calculable par une machine de Turing.
\end{enumerate}
Ces fonctions sont les \defin{fonctions calculables} au sens de
Church-Turing.

Les équivalences sont constructives~: elles fournissent des traductions
algorithmiques entre les différentes représentations.  On peut les
comprendre comme des compilateurs transformant un programme écrit dans
un modèle en un programme équivalent écrit dans un autre.

\thingy Ces modèles ont également les mêmes limitations.  Il n'existe
pas d'algorithme permettant de déterminer si une machine de Turing
donnée s'arrête, ni de déterminer si la réduction d'un
$\lambda$-terme donné atteint une forme normale.  Plus généralement,
le théorème de Rice montre qu'aucune propriété non triviale de la
fonction calculée par un programme n'est décidable à partir du code de
ce programme.

\thingy Un langage de programmation est dit
\defin[Turing-complet (langage)]{Turing-complet} lorsque, après une
idéalisation convenable, il permet de calculer exactement les fonctions
calculables au sens de Church-Turing.  Les langages généralistes usuels
sont Turing-complets~: Python, Java, JavaScript, C, C++, OCaml, Haskell,
Lisp, Perl, Ruby, Smalltalk ou Prolog, par exemple.

Certains langages ou systèmes le sont plus ou moins accidentellement,
par exemple certains fragments ou usages de CSS, \TeX, XSLT ou
\texttt{gm}.  La Turing-complétude ne signifie cependant pas qu'un
langage peut commodément exprimer n'importe quelle opération.  Elle
signifie seulement qu'il peut, en principe et avec les codages
appropriés, calculer toute fonction calculable.  La traduction peut
être extrêmement inefficace, illisible ou malcommode.

Un ordinateur réel ne peut pas calculer davantage qu'une machine de
Turing, sauf à disposer d'une véritable source de hasard qui modifierait
la nature des résultats attendus.  La question pertinente est plutôt
de savoir si un langage donné permet de calculer \emph{autant} qu'une
machine de Turing.

\thingy Enfin, le typage ne garantit pas à lui seul la terminaison.
L'existence d'un type $t$ tel que
\[
t\cong(t\to t)
\]
permet d'encoder dans le type $t$ les termes du $\lambda$-calcul non
typé.  On peut ainsi provoquer une boucle infinie dans un langage
apparemment typé, sans employer explicitement une définition
récursive.

\review{La construction OCaml donnée dans les diapositives utilise un
type récursif pour représenter l'auto-application.  Elle pourra être
reprise comme exemple ou exercice, mais elle n'est pas reproduite ici
afin de respecter la consigne de ne pas intégrer encore les
exercices.}

\thingy Les fonctions calculables peuvent croître extraordinairement
vite.  La fonction d'Ackermann en est un premier exemple, tandis que la
fonction du castor affairé domine toute fonction calculable.  Il est
donc possible de concevoir un programme de taille raisonnable qui
termine théoriquement, mais dont le temps d'exécution ou le résultat
dépasse toute grandeur utilisable en pratique.

La calculabilité distingue ainsi deux questions radicalement
différentes~:
\begin{itemize}
\item un calcul peut-il être effectué par un algorithme~?
\item s'il le peut, peut-il être effectué avec des ressources
  raisonnables~?
\end{itemize}
La première relève de la calculabilité~; la seconde, de la complexité.